\documentclass[11pt]{article}
\usepackage[margin=2.5cm]{geometry}
\usepackage{amsmath,amssymb}
\usepackage[expansion=false]{microtype}
\usepackage[hidelinks]{hyperref}
\usepackage{booktabs}
\usepackage{enumitem}
\usepackage{xcolor}
\usepackage{listings}

\definecolor{added}{rgb}{0.0, 0.45, 0.0}
\definecolor{codebg}{rgb}{0.96, 0.96, 0.96}

\lstset{
  backgroundcolor=\color{codebg},
  basicstyle=\ttfamily\small,
  breaklines=true,
  frame=single,
  rulecolor=\color{black!30},
  xleftmargin=0.5em,
  xrightmargin=0.5em,
  aboveskip=0.8em,
  belowskip=0.8em,
}

\newcommand{\specref}[1]{\textsection#1}
\newcommand{\field}[1]{\texttt{#1}}
\newcommand{\struct}[1]{\texttt{#1}}

\title{\textbf{Principia Spec Delta} \\[6pt]
\large COM Projection, Invariant Monitoring, and \\
Physics Visualisation Modes \\[4pt]
\normalsize Integration into the Implementation Spec}
\author{Malachy Doherty}
\date{March 2026}

\begin{document}
\maketitle

\begin{abstract}
This document specifies the exact changes required to integrate COM projection,
invariant monitoring, and a comprehensive set of physics visualisation modes
into the existing Principia v1 implementation spec. It covers: additions to
the \struct{SampleSummary} struct, modifications to the per-sample integration
kernel, new uniform buffer fields, a full taxonomy of colour mapping modes for
physical quantities (energy, angular momentum, mass ratios, momenta, Jacobi
vectors, and derived diagnostics), and GUI surface changes. Each section
identifies the insertion point in the existing spec by section number.
\end{abstract}

\tableofcontents
\newpage

% ============================================================
\section{Overview of touched spec sections}

\begin{center}
\begin{tabular}{lll}
\toprule
Existing spec section & What changes & Nature of change \\
\midrule
\struct{SampleSummary} struct & Diagnostic + physics fields & Additive \\
Per-sample compute kernel & Projection step + invariant accumulation & Insert \\
Per-sample compute kernel & IC-derived physics field writes & Insert \\
Uniform buffer layout & Stabilising floors $\varepsilon_E$, $\varepsilon_L$ & Additive \\
Colour mapping / overlays & Physics vis modes + drift overlays & Additive \\
Data export format & New fields in binary header & Additive \\
GUI diagnostic panel & Invariant health readout & Additive \\
GUI colour mode selector & Physics visualisation group & Additive \\
\bottomrule
\end{tabular}
\end{center}

No existing fields are removed or reinterpreted. All changes are additive or
insertions into the existing integration loop.

% ============================================================
\section{\texttt{SampleSummary} struct changes}
\label{sec:struct_changes}

\textit{Insertion point: the existing \struct{SampleSummary} struct definition.}

Add the following fields to the per-sample output struct:

\begin{lstlisting}[language=C]
// --- COM projection and invariant diagnostics ---

// Energy drift
f32 delta_E_final;    // E(t_end) - E(0)
f32 delta_E_max_abs;  // max_{tau <= t_end} |E(tau) - E(0)|

// Angular momentum drift
f32 delta_Lz_final;   // Lz(t_end) - Lz(0)
f32 delta_Lz_max_abs; // max_{tau <= t_end} |Lz(tau) - Lz(0)|

// Relative drift (for colour mapping)
f32 delta_E_rel;      // delta_E_final / max(eps_E, |E(0)|)
f32 delta_Lz_rel;     // delta_Lz_final / max(eps_L, |Lz(0)|)

// Initial invariant values (needed for relative drift
// computation and for downstream analysis)
f32 E_0;              // E(0)
f32 Lz_0;             // Lz(0)

// --- IC-derived physics quantities ---
// (computed once from the decoded initial condition,
//  before the integration loop begins)

// Masses
f32 m1;               // mass of body 1
f32 m2;               // mass of body 2
f32 m3;               // mass of body 3

// Mass ratios (chart-dependent; only vary across the
// slice for charts that parameterise mass)
f32 mu_12;            // m1 / m2
f32 mu_13;            // m1 / m3
f32 q_mass;           // m_min / M_total  (mass fraction
                      // of lightest body: 0 = test particle,
                      // 1/3 = equal mass)

// Jacobi vectors (from the decoded IC)
f32 rho1_mag;         // |rho_1|
f32 rho2_mag;         // |rho_2|
f32 rho_ratio;        // |rho_2| / |rho_1|
f32 rho_angle;        // angle between rho_1 and rho_2 [rad]

// Momenta
f32 K_0;              // total kinetic energy at t=0
f32 V_0;              // gravitational potential at t=0
f32 virial_ratio;     // 2*K_0 / |V_0|  (= 1 for virial
                      // equilibrium)

// Pairwise geometry
f32 r_min_pair_0;     // min_{i<j} |r_i - r_j| at t=0
                      // (closest pair distance)
f32 r_min_pair_final; // min_{i<j} |r_i - r_j| over the
                      // whole trajectory (closest approach)

// Trajectory-derived
f32 escape_time;      // time to termination (escape,
                      // collision, or horizon)
u32 n_close;          // number of close encounters
                      // (r_ij < r_close threshold)
\end{lstlisting}

\paragraph{Sizing.} The invariant diagnostics add eight \texttt{f32} fields
(32 bytes). The IC-derived physics quantities add a further fifteen \texttt{f32}
fields and one \texttt{u32} field (64 bytes). Total additional payload: 96 bytes
per sample. For a $1024 \times 1024$ tile this is $\sim$96\,MB, which is modest
relative to the existing \struct{SampleSummary} payload including checkpoint
arrays.

\paragraph{Bit-packing note.} These fields are continuous \texttt{f32} values,
not bit-packed flags. They sit alongside the existing bit-packed
\field{descriptor} and \field{outcome} fields without interfering. If struct
alignment requires it, pad to a 16-byte boundary after the last new field.

\paragraph{Design rationale: why both absolute and relative drift.}
Absolute drift (\field{delta\_E\_final}, \field{delta\_Lz\_final}) is the
physically meaningful quantity. Relative drift (\field{delta\_E\_rel},
\field{delta\_Lz\_rel}) is useful for colour mapping because it normalises
across regions of the IC manifold where $|E(0)|$ varies by orders of magnitude.
The max-abs fields (\field{delta\_E\_max\_abs}, \field{delta\_Lz\_max\_abs})
catch transient close-encounter spikes that may recover by the end of the
trajectory---exactly the regime where the Burrau dynamics are most interesting.

% ============================================================
\section{Uniform buffer additions}
\label{sec:uniforms}

\textit{Insertion point: the existing per-dispatch uniform buffer layout.}

Add to the simulation parameter uniform block:

\begin{lstlisting}[language=C]
// Stabilising floors for relative drift computation.
// Prevents division by near-zero when E(0) or Lz(0)
// is close to zero (e.g. parabolic energy, zero angular
// momentum rest configurations).
f32 eps_E;     // default: 1e-6
f32 eps_L;     // default: 1e-6

// Close-encounter threshold for n_close counting.
f32 r_close;   // default: 0.01 (normalised units)
\end{lstlisting}

These are set once at dispatch time by the CPU and do not change per-step.
Default values are appropriate for the normalised unit system used by the atlas
decoder. If the simulator later supports physical units, these should scale
accordingly.

% ============================================================
\section{Integration kernel modifications}
\label{sec:kernel}

\textit{Insertion point: the per-sample integration loop, between the final
integration step and the \struct{SampleSummary} write-out.}

\subsection{Initialisation (before the integration loop)}

After decoding the initial condition and before entering the timestep loop,
evaluate and store the initial invariants:

\begin{lstlisting}[language=C]
// --- Invariant initialisation ---
let E_0: f32 = compute_energy(r, v, m);  // Eq. (18) of mini spec
let Lz_0: f32 = compute_Lz(r, v, m);    // Eq. (19) of mini spec

// Running max-abs accumulators
var delta_E_max_abs: f32 = 0.0;
var delta_Lz_max_abs: f32 = 0.0;
\end{lstlisting}

Where \texttt{compute\_energy} and \texttt{compute\_Lz} are:

\begin{lstlisting}[language=C]
fn compute_energy(
    r: array<vec2<f32>, 3>,
    v: array<vec2<f32>, 3>,
    m: array<f32, 3>
) -> f32 {
    // Kinetic energy
    var T: f32 = 0.0;
    for (var i = 0u; i < 3u; i++) {
        T += 0.5 * m[i] * dot(v[i], v[i]);
    }

    // Gravitational potential (G = 1 in natural units)
    var V: f32 = 0.0;
    for (var i = 0u; i < 3u; i++) {
        for (var j = i + 1u; j < 3u; j++) {
            let dr = r[i] - r[j];
            let dist = length(dr);
            V -= m[i] * m[j] / dist;
        }
    }

    return T + V;
}

fn compute_Lz(
    r: array<vec2<f32>, 3>,
    v: array<vec2<f32>, 3>,
    m: array<f32, 3>
) -> f32 {
    var Lz: f32 = 0.0;
    for (var i = 0u; i < 3u; i++) {
        // 2D cross product: x * vy - y * vx
        Lz += m[i] * (r[i].x * v[i].y - r[i].y * v[i].x);
    }
    return Lz;
}
\end{lstlisting}

\paragraph{Note on force reuse.} The pairwise distances $\|r_i - r_j\|$ and
the mass products $m_i m_j$ computed here are the same quantities needed for the
gravitational force evaluation. If the integrator caches these (which it should
for the leapfrog/symplectic step), the potential energy evaluation is
essentially free---just sum the cached $-m_i m_j / \|r_i - r_j\|$ terms. The
kinetic energy term requires three dot products, which is negligible.

\subsection{IC-derived physics quantities (before the integration loop)}

Immediately after invariant initialisation, compute and store the per-sample
physics quantities that derive from the decoded initial condition. These are
evaluated once and written directly to \struct{SampleSummary}; they do not
change during integration.

\begin{lstlisting}[language=C]
// --- IC-derived physics quantities ---
// (r, v, m are the decoded initial state;
//  rho1, rho2 are the Jacobi vectors from the decoder)

// Masses (store for downstream use; the decoder
// already has these)
summary.m1 = m[0];
summary.m2 = m[1];
summary.m3 = m[2];

// Mass ratios
summary.mu_12 = m[0] / m[1];
summary.mu_13 = m[0] / m[2];
let m_min = min(m[0], min(m[1], m[2]));
summary.q_mass = m_min / M_total;

// Jacobi vectors
// rho1 and rho2 come from the atlas decoder;
// if the integrator works in Cartesian coordinates,
// reconstruct from positions:
//   rho1 = r[1] - r[0]
//   com12 = (m[0]*r[0] + m[1]*r[1]) / (m[0] + m[1])
//   rho2 = r[2] - com12
let rho1 = r[1] - r[0];
let com12 = (m[0] * r[0] + m[1] * r[1]) / (m[0] + m[1]);
let rho2 = r[2] - com12;

summary.rho1_mag = length(rho1);
summary.rho2_mag = length(rho2);
summary.rho_ratio = summary.rho2_mag
    / max(1e-30, summary.rho1_mag);
// Signed angle from rho1 to rho2
summary.rho_angle = atan2(
    rho1.x * rho2.y - rho1.y * rho2.x,
    dot(rho1, rho2)
);

// Energetics
// (E_0 already computed in invariant initialisation)
var K_0: f32 = 0.0;
for (var i = 0u; i < 3u; i++) {
    K_0 += 0.5 * m[i] * dot(v[i], v[i]);
}
summary.K_0 = K_0;
summary.V_0 = E_0 - K_0;  // V = E - T
summary.virial_ratio = 2.0 * K_0
    / max(1e-30, abs(summary.V_0));

// Pairwise geometry at t=0
var r_min_0: f32 = 1e30;
for (var i = 0u; i < 3u; i++) {
    for (var j = i + 1u; j < 3u; j++) {
        let d = length(r[i] - r[j]);
        r_min_0 = min(r_min_0, d);
    }
}
summary.r_min_pair_0 = r_min_0;

// Trajectory-derived accumulators (initialise here,
// update during integration)
var r_min_pair_running: f32 = r_min_0;
var n_close: u32 = 0u;
\end{lstlisting}

\paragraph{Jacobi vector convention.} The Jacobi decomposition pairs bodies
$1$--$2$ as the inner pair and body $3$ as the outer, consistent with the atlas
decoder convention. If the decoder uses a different pairing (e.g.\ for
different chart families), the same formulae apply with the appropriate index
permutation.

\paragraph{Close-encounter counting.} The threshold $r_{\text{close}}$ for
incrementing \field{n\_close} should be a uniform parameter (default:
$r_{\text{close}} = 0.01$ in normalised units). Inside the integration loop,
after each step:

\begin{lstlisting}[language=C]
// --- Close encounter tracking (inside loop) ---
for (var i = 0u; i < 3u; i++) {
    for (var j = i + 1u; j < 3u; j++) {
        let d = length(r[i] - r[j]);
        r_min_pair_running = min(r_min_pair_running, d);
        if (d < uniforms.r_close) {
            n_close += 1u;
        }
    }
}
\end{lstlisting}

And at trajectory termination:
\begin{lstlisting}[language=C]
summary.r_min_pair_final = r_min_pair_running;
summary.n_close = n_close;
summary.escape_time = t;  // current simulation time
\end{lstlisting}

\subsection{COM projection (inside the integration loop)}

After each integration substep, insert:

\begin{lstlisting}[language=C]
// --- COM projection (every timestep) ---
// Position: remove COM drift
var R_com = vec2<f32>(0.0);
var M_total: f32 = 0.0;
for (var i = 0u; i < 3u; i++) {
    R_com += m[i] * r[i];
    M_total += m[i];
}
R_com /= M_total;

for (var i = 0u; i < 3u; i++) {
    r[i] -= R_com;
}

// Velocity: remove bulk momentum drift
var V_com = vec2<f32>(0.0);
for (var i = 0u; i < 3u; i++) {
    V_com += m[i] * v[i];
}
V_com /= M_total;

for (var i = 0u; i < 3u; i++) {
    v[i] -= V_com;
}
\end{lstlisting}

\paragraph{Integrator placement.} For a leapfrog (Stormer--Verlet) integrator
with the kick-drift-kick structure, the projection should occur after the full
step is complete (i.e.\ after the second half-kick), not between the half-kick
and the drift. This ensures the symplectic structure of the individual step is
not broken. The projection itself is a canonical linear map, so it commutes with
the symplectic structure in the exact-arithmetic limit; in \texttt{f32} the
ordering matters only at the roundoff level, but placing it after the full step
is the cleanest convention.

\paragraph{Total mass caching.} $M_{\text{total}}$ is constant across the
trajectory. Compute it once during initialisation and pass it as a local
variable rather than recomputing per step.

\subsection{Invariant accumulation (inside the integration loop)}

After the COM projection and before the next timestep, evaluate the current
invariant drift and update the running maximum:

\begin{lstlisting}[language=C]
// --- Invariant accumulation (every timestep) ---
let E_now: f32 = compute_energy(r, v, m);
let Lz_now: f32 = compute_Lz(r, v, m);

let dE = abs(E_now - E_0);
let dLz = abs(Lz_now - Lz_0);

delta_E_max_abs = max(delta_E_max_abs, dE);
delta_Lz_max_abs = max(delta_Lz_max_abs, dLz);
\end{lstlisting}

\paragraph{Cost.} One energy evaluation and one $L_z$ evaluation per timestep.
As noted above, the pairwise distances are already available from the force
computation, so the marginal cost is roughly six multiplies, three additions,
three dot products, and two comparisons per step. This is negligible relative
to the force evaluation and integration arithmetic.

\subsection{Final write-out (after the integration loop)}

After the trajectory terminates (escape, collision, or horizon reached),
compute the final drift values and write to \struct{SampleSummary}:

\begin{lstlisting}[language=C]
// --- Final invariant diagnostics ---
let E_final = compute_energy(r, v, m);
let Lz_final = compute_Lz(r, v, m);

summary.E_0 = E_0;
summary.Lz_0 = Lz_0;

summary.delta_E_final = E_final - E_0;
summary.delta_Lz_final = Lz_final - Lz_0;

summary.delta_E_max_abs = delta_E_max_abs;
summary.delta_Lz_max_abs = delta_Lz_max_abs;

// Relative drift with stabilising floor
summary.delta_E_rel =
    summary.delta_E_final / max(uniforms.eps_E, abs(E_0));
summary.delta_Lz_rel =
    summary.delta_Lz_final / max(uniforms.eps_L, abs(Lz_0));
\end{lstlisting}

% ============================================================
\section{Physics visualisation modes}
\label{sec:colour}

\textit{Insertion point: the existing colour mapping specification, alongside
the physics overlay modes.}

The existing colour mapping system supports basin colouring, shape sphere,
and angular velocity differential modes. This section adds a comprehensive
set of physics visualisation modes organised into four groups. Each mode maps
a per-sample scalar from \struct{SampleSummary} to a colour via a 1D palette
lookup.

\subsection{Group 1: Conserved quantities and invariants}

These modes visualise the physical invariants of each initial condition. They
reveal the energy and angular-momentum landscape of the IC manifold---structure
that is invisible in the basin colouring but physically fundamental.

\begin{center}
\begin{tabular}{llll}
\toprule
Mode name & Source field & Palette & Notes \\
\midrule
\texttt{ENERGY} & \field{E\_0} & Diverging & Negative bound,
    positive unbound. \\
 & & (centred at $E=0$) & Zero contour = parabolic. \\
\texttt{ANG\_MOMENTUM} & \field{Lz\_0} & Diverging & Sign distinguishes \\
 & & (centred at $L_z=0$) & prograde/retrograde. \\
\texttt{KINETIC} & \field{K\_0} & Sequential & Always $\geq 0$. \\
\texttt{POTENTIAL} & \field{V\_0} & Sequential & Always $\leq 0$. \\
 & & (inverted) & More negative = deeper well. \\
\texttt{VIRIAL} & \field{virial\_ratio} & Diverging & $= 1$ at virial equilibrium. \\
 & & (centred at 1) & $< 1$ bound-dominated. \\
\bottomrule
\end{tabular}
\end{center}

\paragraph{Energy colour scale.} The energy range across a typical Burrau-family
slice spans several orders of magnitude (deeply bound binaries to marginally
unbound escapers). A symmetric log scale is recommended: linear near zero
(within $[-\varepsilon, +\varepsilon]$), logarithmic outside. The zero-energy
contour should be visually distinct (e.g.\ a thin white or black line) since it
separates bound and unbound configurations.

\paragraph{Angular momentum colour scale.} $L_z$ is signed. A diverging palette
centred at zero naturally distinguishes prograde ($L_z > 0$) from retrograde
($L_z < 0$) configurations. For the Burrau rest configurations $L_z = 0$
everywhere, so this mode becomes informative only when the chart parameterises
initial momentum.

\subsection{Group 2: Mass and configuration geometry}

These modes visualise the shape and mass structure of the initial configuration.
They are most informative on charts that vary mass or geometry across the slice
(e.g.\ the mass-simplex chart, the Burrau $\nu$-sweep).

\begin{center}
\begin{tabular}{llll}
\toprule
Mode name & Source field & Palette & Notes \\
\midrule
\texttt{MASS\_RATIO\_12} & \field{mu\_12} & Diverging & Centred at 1 \\
 & & (centred at 1) & (equal mass pair). \\
\texttt{MASS\_RATIO\_13} & \field{mu\_13} & Diverging & Centred at 1. \\
\texttt{MASS\_FRACTION} & \field{q\_mass} & Sequential & 0 = test particle, \\
 & & & $1/3$ = equal mass. \\
\texttt{JACOBI\_RHO1} & \field{rho1\_mag} & Sequential & Inner pair separation. \\
\texttt{JACOBI\_RHO2} & \field{rho2\_mag} & Sequential & Outer body distance. \\
\texttt{JACOBI\_RATIO} & \field{rho\_ratio} & Log-scale & $|\rho_2|/|\rho_1|$. \\
 & & diverging & $> 1$: hierarchical. \\
 & & (centred at 1) & $< 1$: compact. \\
\texttt{JACOBI\_ANGLE} & \field{rho\_angle} & Cyclic & $[-\pi, \pi]$. \\
 & & & Use a cyclic palette \\
 & & & (e.g.\ HSL hue wheel). \\
\texttt{MIN\_PAIR\_DIST} & \field{r\_min\_pair\_0} & Sequential & Closest pair at $t=0$. \\
 & & (log scale) & Small = near-collision IC. \\
\bottomrule
\end{tabular}
\end{center}

\paragraph{Jacobi ratio interpretation.} The ratio $|\rho_2|/|\rho_1|$ is a
scale-free measure of how hierarchical the configuration is. Values $\gg 1$
indicate a tight inner pair with a distant third body (the hierarchical triple
regime where perturbation theory works well). Values $\sim 1$ indicate a compact
configuration where all three bodies interact comparably (the chaotic regime).
This is arguably the single most physically informative scalar you can extract
from the geometry.

\paragraph{Jacobi angle colour scale.} Since the angle wraps at $\pm\pi$, the
palette must be cyclic---the colour at $-\pi$ and $+\pi$ must match. Standard
cyclic palettes (twilight, HSL hue, phase colourmaps) work here. This mode
should integrate with the existing colour mapping system's OKLAB pipeline; a
cyclic VMF-OKLAB mode would be the natural implementation.

\subsection{Group 3: Trajectory-derived quantities}

These modes visualise properties of the integrated trajectory, not just the
initial condition. They require the integration to have completed.

\begin{center}
\begin{tabular}{llll}
\toprule
Mode name & Source field & Palette & Notes \\
\midrule
\texttt{ESCAPE\_TIME} & \field{escape\_time} & Sequential & Time to termination. \\
 & & (log scale) & Long-lived = fractal boundary. \\
\texttt{CLOSE\_ENCOUNTERS} & \field{n\_close} & Sequential & Count of close approaches. \\
 & & (integer) & High = chaotic scattering. \\
\texttt{MIN\_APPROACH} & \field{r\_min\_pair\_final} & Sequential & Closest approach over \\
 & & (log scale) & entire trajectory. \\
\bottomrule
\end{tabular}
\end{center}

\paragraph{Escape time and fractal dimension.} The escape time map is closely
related to the fractal dimension of basin boundaries. Pixels with long escape
times cluster on the boundaries between basins---they are the orbits that take
longest to ``decide'' which outcome to reach. Rendering escape time in log scale
with a perceptually uniform palette produces striking images where the fractal
structure of the boundary is visible as luminance variation.

\paragraph{Close encounter count.} This is an integer field. The palette should
use discrete steps for small counts ($0$--$10$) and transition to a continuous
gradient for large counts. Regions with high encounter counts correlate with
chaotic scattering and long-lived triple interactions.

\subsection{Group 4: Numerical diagnostics}

These modes visualise the numerical fidelity of the integration, as specified
in the COM projection mini spec.

\begin{center}
\begin{tabular}{llll}
\toprule
Mode name & Source field & Palette & Notes \\
\midrule
\texttt{ENERGY\_DRIFT\_ABS} & \field{delta\_E\_max\_abs} & Log-scale & Absolute worst-case. \\
\texttt{ENERGY\_DRIFT\_REL} & \field{delta\_E\_rel} & Log-scale & Normalised by $|E(0)|$. \\
\texttt{LZ\_DRIFT\_ABS} & \field{delta\_Lz\_max\_abs} & Log-scale & Absolute worst-case. \\
\texttt{LZ\_DRIFT\_REL} & \field{delta\_Lz\_rel} & Log-scale & Normalised by $|L_z(0)|$. \\
\bottomrule
\end{tabular}
\end{center}

\paragraph{Diagnostic colour scale.} A sequential log-scale palette:
\begin{itemize}[nosep]
    \item Floor at $\sim 10^{-7}$ (single-precision machine epsilon regime).
    \item Ceiling at $\sim 10^{0}$ (drift comparable to the invariant itself).
    \item Regions with drift $> 10^{-2}$ should be visually distinct
          (warm/hot colours) as these indicate numerically suspect trajectories.
\end{itemize}

\subsection{Compositing model}

All physics visualisation modes operate in two regimes:

\paragraph{Primary mode.} The mode replaces the basin colouring entirely. The
palette maps the scalar field to RGB directly. This is useful when the physics
quantity itself is the object of study (e.g.\ ``show me the energy landscape'').

\paragraph{Overlay mode.} The mode is composited over the existing primary
render (basin, shape sphere, etc.) as a semi-transparent layer. The alpha
channel is derived from the scalar magnitude: regions where the quantity is
near its midpoint or zero are transparent, regions where it is extreme are
opaque. This is useful for ``show me the basin map, but highlight where the
energy is extreme'' workflows.

The compositing pipeline should use premultiplied alpha and operate in linear
OKLAB space to avoid hue shifts at intermediate opacities.

\subsection{GUI integration}

The colour mode selector should organise the new modes into named groups:

\begin{center}
\begin{tabular}{ll}
\toprule
Group & Modes \\
\midrule
Invariants & \texttt{ENERGY}, \texttt{ANG\_MOMENTUM}, \texttt{KINETIC}, \\
 & \texttt{POTENTIAL}, \texttt{VIRIAL} \\
Geometry & \texttt{MASS\_RATIO\_12}, \texttt{MASS\_RATIO\_13}, \\
 & \texttt{MASS\_FRACTION}, \texttt{JACOBI\_RHO1}, \\
 & \texttt{JACOBI\_RHO2}, \texttt{JACOBI\_RATIO}, \\
 & \texttt{JACOBI\_ANGLE}, \texttt{MIN\_PAIR\_DIST} \\
Trajectory & \texttt{ESCAPE\_TIME}, \texttt{CLOSE\_ENCOUNTERS}, \\
 & \texttt{MIN\_APPROACH} \\
Diagnostics & \texttt{ENERGY\_DRIFT\_ABS}, \texttt{ENERGY\_DRIFT\_REL}, \\
 & \texttt{LZ\_DRIFT\_ABS}, \texttt{LZ\_DRIFT\_REL} \\
\bottomrule
\end{tabular}
\end{center}

Each group is collapsible in the UI. The ``Diagnostics'' group should be
selectable as an overlay independently of the other groups (i.e.\ you can have
\texttt{ENERGY} as primary + \texttt{ENERGY\_DRIFT\_REL} as overlay
simultaneously). The other groups are mutually exclusive with the existing
primary modes (basin, shape sphere, angular velocity differential).

% ============================================================
\section{GUI diagnostic panel}
\label{sec:gui}

\textit{Insertion point: the GUI spec, as a new panel or section of the
existing debug/info panel.}

\subsection{Global health indicators}

After a tile or frame completes rendering, the CPU-side reduction should compute
and display:

\begin{center}
\begin{tabular}{ll}
\toprule
Indicator & Computation \\
\midrule
Worst energy drift & $\max_{\text{samples}} |\field{delta\_E\_max\_abs}|$ \\
Worst $L_z$ drift & $\max_{\text{samples}} |\field{delta\_Lz\_max\_abs}|$ \\
Median energy drift & median of $|\field{delta\_E\_max\_abs}|$ over samples \\
Median $L_z$ drift & median of $|\field{delta\_Lz\_max\_abs}|$ over samples \\
Suspect sample fraction & fraction with $|\field{delta\_E\_rel}| > 10^{-2}$ \\
\bottomrule
\end{tabular}
\end{center}

\paragraph{Display.} These belong in a collapsible ``Numerical health'' panel,
visible by default in research mode, hidden by default in presentation mode.
The worst-case values should have traffic-light colouring: green below
$10^{-5}$, amber between $10^{-5}$ and $10^{-2}$, red above $10^{-2}$.

\subsection{Per-pixel tooltip}

When the user hovers over a rendered pixel (if the cursor inspector is active),
the tooltip should additionally display:

\begin{itemize}[nosep]
    \item $E(0)$, $L_z(0)$, $K(0)$, $V(0)$ for that sample.
    \item Mass triple $(m_1, m_2, m_3)$ and mass fraction $q$.
    \item Jacobi magnitudes $|\rho_1|$, $|\rho_2|$, ratio, and angle.
    \item Virial ratio $2K/|V|$.
    \item Escape time and close-encounter count.
    \item $\Delta E_{\max}$ and $\Delta L_{z,\max}$ (the transient worst-case drift).
    \item $\delta E_{\text{rel}}$ and $\delta L_{\text{rel}}$ (the relative final drift).
\end{itemize}

This is a lot of data. The tooltip should be organised into collapsible groups
matching the four colour mode groups (Invariants, Geometry, Trajectory,
Diagnostics), with only the currently active group expanded by default.

% ============================================================
\section{Data export format changes}
\label{sec:export}

\textit{Insertion point: the binary export format specification (\specref{26.2.4}).}

\subsection{Binary header}

Add to the export metadata header:

\begin{lstlisting}[language=C]
f32 eps_E;             // stabilising floor used
f32 eps_L;             // stabilising floor used
f32 r_close;           // close-encounter threshold used
u32 com_projection;    // 1 = projection active, 0 = disabled
u32 invariant_cadence; // steps between invariant evaluations
                       // (1 = every step, default)
\end{lstlisting}

\subsection{Per-sample record}

All 24 new \struct{SampleSummary} fields are appended to each per-sample
record in the binary export. The field order matches the struct definition in
\S\ref{sec:struct_changes}. Downstream consumers (e.g.\ \texttt{np.frombuffer})
need the updated struct dtype to parse correctly.

A recommended NumPy dtype fragment for the new fields:

\begin{lstlisting}[language=Python]
physics_dtype = np.dtype([
    # Invariant diagnostics
    ('delta_E_final', 'f4'), ('delta_E_max_abs', 'f4'),
    ('delta_Lz_final', 'f4'), ('delta_Lz_max_abs', 'f4'),
    ('delta_E_rel', 'f4'), ('delta_Lz_rel', 'f4'),
    ('E_0', 'f4'), ('Lz_0', 'f4'),
    # Masses
    ('m1', 'f4'), ('m2', 'f4'), ('m3', 'f4'),
    # Mass ratios
    ('mu_12', 'f4'), ('mu_13', 'f4'), ('q_mass', 'f4'),
    # Jacobi vectors
    ('rho1_mag', 'f4'), ('rho2_mag', 'f4'),
    ('rho_ratio', 'f4'), ('rho_angle', 'f4'),
    # Energetics
    ('K_0', 'f4'), ('V_0', 'f4'), ('virial_ratio', 'f4'),
    # Pairwise geometry
    ('r_min_pair_0', 'f4'), ('r_min_pair_final', 'f4'),
    # Trajectory-derived
    ('escape_time', 'f4'), ('n_close', 'u4'),
])
\end{lstlisting}

\subsection{CSV / JSON export}

For text-format exports, add all new fields as columns/keys. Group them
with a naming prefix for clarity:

\begin{lstlisting}
# Invariant diagnostics
E_0, Lz_0, delta_E_final, delta_Lz_final,
delta_E_max_abs, delta_Lz_max_abs,
delta_E_rel, delta_Lz_rel

# Mass and geometry
m1, m2, m3, mu_12, mu_13, q_mass,
rho1_mag, rho2_mag, rho_ratio, rho_angle,
r_min_pair_0

# Energetics
K_0, V_0, virial_ratio

# Trajectory-derived
escape_time, n_close, r_min_pair_final
\end{lstlisting}

% ============================================================
\section{Implementation plan integration}
\label{sec:plan}

\textit{Where this lands in the milestone-ordered task breakdown.}

The COM projection, invariant monitoring, and physics visualisation work splits
across existing milestones:

\begin{center}
\begin{tabular}{lll}
\toprule
Task & Milestone & Size \\
\midrule
\multicolumn{3}{l}{\textit{Compute pipeline (M1)}} \\
Add diagnostic fields to \struct{SampleSummary} & M1 & S \\
Add IC-derived physics fields to \struct{SampleSummary} & M1 & S \\
Add trajectory-derived fields to \struct{SampleSummary} & M1 & S \\
Implement \texttt{compute\_energy}, \texttt{compute\_Lz} & M1 & S \\
Insert COM projection into integration loop & M1 & S \\
Insert invariant accumulation into loop & M1 & S \\
Compute IC-derived physics quantities after decode & M1 & S \\
Close-encounter counting in integration loop & M1 & S \\
Wire all new fields to \struct{SampleSummary} write-out & M1 & S \\
Add $\varepsilon_E$, $\varepsilon_L$, $r_{\text{close}}$ to uniforms & M1 & S \\
\midrule
\multicolumn{3}{l}{\textit{Colour mapping (M3)}} \\
Implement Group 1 modes (invariants) & M3 & M \\
Implement Group 2 modes (geometry) & M3 & M \\
Implement Group 3 modes (trajectory) & M3 & S \\
Implement Group 4 modes (diagnostics) & M3 & S \\
Cyclic palette for Jacobi angle mode & M3 & S \\
Symmetric log scale for energy mode & M3 & S \\
Primary/overlay compositing pipeline & M3 & M \\
Colour mode group selector in GUI & M3 & S \\
\midrule
\multicolumn{3}{l}{\textit{GUI (M4)}} \\
CPU-side reduction for global health indicators & M4 & S \\
Numerical health panel in GUI & M4 & S \\
Per-pixel tooltip: all physics fields, grouped & M4 & M \\
\midrule
\multicolumn{3}{l}{\textit{Export (M5)}} \\
Update binary export header and per-sample record & M5 & S \\
Update CSV/JSON export with all new fields & M5 & S \\
\bottomrule
\end{tabular}
\end{center}

The M1 compute tasks should be done as a single unit of work---they are tightly
coupled in the shader. The M3 colour mapping tasks are the largest block; Groups
1--2 are medium because they require new palette types (symmetric log, cyclic)
and the compositing pipeline. Groups 3--4 reuse existing sequential palettes
and are straightforward. The M4 and M5 tasks can be done independently whenever
those milestones are reached.

% ============================================================
\section{Validation}
\label{sec:validation}

To confirm the implementation is correct:

\begin{enumerate}[nosep]
    \item \textbf{COM conservation.} After projection, $\|R_{\text{com}}\|$ and
    $\|P_{\text{com}}\|$ should be at or below \texttt{f32} machine epsilon
    ($\sim 10^{-7}$) times the characteristic scale. If not, the projection
    arithmetic is wrong.

    \item \textbf{Known orbit test.} For the Lagrange equilateral triangle
    solution (equal masses, circular orbit), $E$ and $L_z$ are exactly conserved
    analytically. The numerical drift should be bounded and oscillatory, not
    secularly growing. Run for $\sim 10^4$ steps and verify
    $|\Delta E| / |E(0)| < 10^{-4}$ (for a symplectic integrator in \texttt{f32}
    this is a reasonable expectation).

    \item \textbf{Euler solution test.} For the collinear Euler solution,
    $L_z = 0$ exactly. Verify that $L_z(0) \approx 0$ and that $\delta L$
    uses the stabilising floor correctly without producing $\text{NaN}$ or
    $\pm\infty$.

    \item \textbf{Close encounter spike.} For a Burrau-family IC known to
    produce a close triple encounter, verify that
    $\field{delta\_E\_max\_abs} > \field{|delta\_E\_final|}$---i.e.\ the
    transient spike is larger than the final drift, confirming the per-step
    accumulator is working.

    \item \textbf{Jacobi vector consistency.} For a known configuration (e.g.\
    the 3--4--5 Burrau triangle with unit hypotenuse), verify that the stored
    Jacobi magnitudes, ratio, and angle match the analytic values from the
    Euclid parametrisation formulae. The angle between $\rho_1$ and $\rho_2$
    for the right-angle vertex at the origin should be exactly
    $\pi/2 + \arctan(a/b)$ (or its equivalent), which is straightforward to
    check.

    \item \textbf{Virial ratio for rest configurations.} For any Burrau-family
    IC with zero initial velocity, $K(0) = 0$ and therefore the virial ratio
    should be exactly zero. Verify that \field{virial\_ratio} $= 0$ (not
    \texttt{NaN} or \texttt{Inf}) when $K(0) = 0$, confirming the
    $\max(10^{-30}, |V(0)|)$ floor in the denominator works correctly.

    \item \textbf{Close encounter count.} For the classic 3--4--5 Burrau
    problem, verify that \field{n\_close} $> 0$ (it is known to produce
    close triple encounters). For a hierarchical triple with well-separated
    inner and outer orbits, verify \field{n\_close} $= 0$ at reasonable
    threshold values.

    \item \textbf{Export round-trip.} Export a tile in binary format, reload in
    Python with the updated dtype, and verify all fields (diagnostic and
    physics) are present, correctly aligned, and match the on-screen tooltip
    values.
\end{enumerate}

% ============================================================
\section{Spec summary}
\label{sec:summary}

The following additions are made to the Principia v1 implementation spec:

\begin{itemize}[nosep]
    \item \struct{SampleSummary} gains 24 new fields (8 invariant diagnostics,
    13 IC-derived physics quantities, 3 trajectory-derived quantities),
    totalling 96 bytes per sample.
    \item The uniform buffer gains two \texttt{f32} stabilising floors
    ($\varepsilon_E$, $\varepsilon_L$), one close-encounter threshold
    ($r_{\text{close}}$), and two \texttt{u32} configuration flags.
    \item The per-sample integration kernel gains: a COM projection step
    (after each full integration step), an invariant accumulation step,
    IC-derived physics computation (after decode, before integration),
    and close-encounter tracking (inside the integration loop).
    \item Two helper functions (\texttt{compute\_energy},
    \texttt{compute\_Lz}) are added to the shader library.
    \item Twenty new colour visualisation modes are added, organised into four
    groups: Invariants (5), Geometry (8), Trajectory (3), and
    Diagnostics (4). These support both primary and overlay rendering with
    compositing in linear OKLAB space.
    \item The GUI gains a ``Numerical health'' panel, a grouped colour mode
    selector, and extended per-pixel tooltip with collapsible physics groups.
    \item The binary, CSV, and JSON export formats gain all new fields.
\end{itemize}

No existing behaviour is modified. The runtime cost is modest: COM projection
adds $\sim 20$ FLOPs per step, the invariant evaluation reuses cached pairwise
distances, and the IC-derived physics quantities are computed once per sample
before integration begins. The close-encounter check adds three distance
comparisons per step, which is negligible.

\end{document}
